\section{The result}\label{sec:result}

How can a model on a lattice give a quantum field theory on continuous space and time?
We study this question for four-dimensional Yang--Mills theory, starting with Wilson's lattice model.
On a finite lattice the variables are finitely many unitary matrices, and their law is an ordinary probability measure.
The coupling is the parameter controlling the weight assigned to the lattice energy; a local observable is a function of matrices near one point.
The problem is to choose them so that their correlations---the averages of their products---have a limit as the lattice becomes finer and the box becomes larger.

We construct such limits for matrix groups with simple nonabelian Lie algebra (Definition~\ref{def:gauge}).
The spacing tends to zero first, at each fixed box size.
From any increasing sequence of box sizes, one can then select a subsequence along which all the required correlations converge.
The limiting correlations are averages of component correlations, called phases.
One selected phase gives quantum fields on a Hilbert space, with positive energy, relativistic covariance and commutativity at spacelike separation (Definition~\ref{def:wightman}).
The same local product expansions, replacing nearby fields by sums of fields at one point, hold in the average and in every phase (Definition~\ref{def:ope}).
A logarithmic correction to the correlation of two copies of one field distinguishes it from the finite free-field expressions specified below.
The result does not establish a mass gap, a positive lower bound on energies above the vacuum.

Part~\ref{part:main} begins with definitions of the lattice model and its continuum objects, followed by the precise theorem.
The strategy then separates the proof into five results: constructing the finite-volume limits, passing to infinite volume and its phases, reconstructing relativistic fields, expanding local products, and proving the logarithmic two-point law.
Part~\ref{part:construction} develops their proofs.
It begins with local lattice observables, studies what happens when short-distance variables are integrated out, and then follows the resulting correlations into the continuum theory.
This order keeps the finite probability model in view while explaining how its limiting correlations become quantum fields.

The statements are matched to the fixed formal-source revision in~\cite{OpenAIYM}.
